\documentclass[12pt,a4paper]{article}
% Сообщение Анастасии: Анастасия, в работе верно выбрана общая идея задач на совместную работу, но во второй задаче перепутано соотношение производительностей; в третьей, пятой и шестой возникли ошибки при преобразовании уравнений; в четвёртой не учтено, что одна труба откачивает воду; в седьмой не прибавлены первые пять дней работы; в восьмой и девятой складывались времена вместо производительностей.
\usepackage[utf8]{inputenc}
\usepackage[T2A]{fontenc}
\usepackage[russian]{babel}
\usepackage[a4paper,margin=2cm]{geometry}
\usepackage{amsmath,amssymb,mathtools}
\usepackage{array,booktabs,longtable}
\usepackage{xcolor}
\usepackage{hyperref}
\usepackage{tikz}
\hypersetup{colorlinks=true,linkcolor=blue!55!black,urlcolor=blue!55!black}
\setlength{\parindent}{0pt}
\setlength{\parskip}{0.55em}
\newcommand{\err}{\textcolor{red}{\textbf{Ошибка:}}}
\newcommand{\idea}{\textcolor{blue!60!black}{\textbf{Ключевая идея:}}}
\newcommand{\result}{\textcolor{teal!60!black}{\textbf{Итог:}}}
\newcommand{\tasksection}[1]{\section*{Задание №#1}}
\begin{document}

\begin{titlepage}
\centering
\vspace*{0.23\textheight}
{\LARGE\bfseries Ревью домашней работы\par}
\vspace{1.8cm}
{\large Автор: Лёвин Артём Александрович\par}
\vspace{0.8cm}
{\large 7 октября 2026 года\par}
\vfill
\begin{tikzpicture}[x=1cm,y=1cm]
\draw[blue!35!black,line width=0.7pt] (-5,0) .. controls (-2.5,0.7) and (2.5,-0.7) .. (5,0);
\end{tikzpicture}
\end{titlepage}

\newpage
\section*{Сводная таблица}
\small
\renewcommand{\arraystretch}{1.18}
\setlength{\tabcolsep}{2pt}
\begin{longtable}{>{\centering\arraybackslash}p{0.08\linewidth} p{0.18\linewidth} p{0.30\linewidth} p{0.18\linewidth} p{0.18\linewidth}}
\toprule
№ & Тема & Суть ошибки & Ваш ответ & Правильный ответ \\
\midrule
\endfirsthead
\toprule
№ & Тема & Суть ошибки & Ваш ответ & Правильный ответ \\
\midrule
\endhead
\hyperref[task:2]{2} & Совместная работа насосов & Общая производительность ошибочно прибавлена к производительности первого насоса вместо вычитания. & $3{,}75$ ч & $15$ ч \\
\midrule
\hyperref[task:3]{3} & Совместная работа насосов & После верно составленного уравнения неверно выполнено избавление от знаменателей; потерян числитель суммы. & В решении получено $2$ ч & $8$ ч и $12$ ч \\
\midrule
\hyperref[task:4]{4} & Наполнение и откачивание & Производительности складывались, хотя вторая труба уменьшает объём воды. & $3$ ч & $6$ ч \\
\midrule
\hyperref[task:5]{5} & Работа станков & При раскрытии скобок и переносе членов изменены знаки; из верного уравнения получено неверное квадратное. & {\scriptsize$\frac{\sqrt{505}-5}{2}$} & $20$ и $15$ деталей/ч \\
\midrule
\hyperref[task:6]{6} & Совместная работа труб & Верное уравнение с дробями преобразовано как линейное: знаменатели фактически потеряны. & $\dfrac{29}{12}$ ч & $10$ ч и $15$ ч \\
\midrule
\hyperref[task:7]{7} & Работа бригад & Найдены $10$ дней после перехода рабочих, но не прибавлены первые $5$ дней. & $10$ дней & $15$ дней \\
\midrule
\hyperref[task:8]{8} & Совместная работа трёх исполнителей & Складывались времена парной работы вместо их производительностей. & Видимый вывод: $14$ ч & $6$ ч \\
\midrule
\hyperref[task:9]{9} & Парная работа трёх мастеров & Из времён парной работы пытались получить индивидуальные времена сложением и вычитанием; нужно работать с обратными величинами. & В таблице: $1{,}75$, $4{,}25$, $5{,}25$ ч & $10$, $15$, $30$ ч \\
\bottomrule
\end{longtable}
\setlength{\tabcolsep}{6pt}
\normalsize

\newpage
\vspace*{0.38\textheight}
\begin{center}
{\LARGE\bfseries Разбор ошибок}
\end{center}
\newpage

\phantomsection\label{task:2}
\tasksection{2}
\subsection*{Текст задания}
Два насоса вместе заполняют резервуар за $6$ часов. Первый насос в одиночку заполняет его за $10$ часов. За сколько часов заполнит резервуар второй насос, работая один?

\subsection*{Ваше решение}
В таблице верно записаны производительности: у первого насоса $\frac{1}{10}$ резервуара в час, у второго $\frac{1}{x}$, а у двух насосов вместе $\frac{1}{6}$. Затем записано
\[
\frac{1}{x}=\frac{1}{10}+\frac{1}{6},
\]
после чего получено $x=3{,}75$.

\subsection*{Ваш ответ}
$3{,}75$ часа.

\subsection*{Ошибка}
\err{} неверно составлено уравнение по таблице производительностей.

\subsection*{Где именно возникает ошибка}
Последний правильный шаг --- таблица с производительностями $\frac{1}{10}$, $\frac{1}{x}$ и $\frac{1}{6}$. Первый неправильный шаг:
\[
\frac{1}{x}=\frac{1}{10}+\frac{1}{6}.
\]

\subsection*{Почему это неверно}
Когда два насоса одновременно наполняют один резервуар, их производительности складываются. Поэтому сумма производительности первого и второго насосов должна быть равна общей производительности. Ваша запись делает производительность второго насоса равной сумме производительности первого насоса и уже общей производительности. Это противоречит смыслу задачи: общая производительность должна быть больше каждой отдельной, а не входить в сумму как ещё один независимый насос.

Здесь известны общая производительность и производительность первого насоса. Значит, производительность второго насоса нужно получить вычитанием первой из общей.

\subsection*{Ключевая идея}
\idea{} для совместной работы складывают именно производительности, а не времена:
\[
\frac{1}{10}+\frac{1}{x}=\frac{1}{6}.
\]

\subsection*{Исправление}
Исправляем только первый неверный шаг:
\[
\frac{1}{10}+\frac{1}{x}=\frac{1}{6}.
\]
Тогда
\[
\frac{1}{x}=\frac{1}{6}-\frac{1}{10}=\frac{5-3}{30}=\frac{2}{30}=\frac{1}{15}.
\]
Следовательно,
\[
x=15.
\]

\subsection*{Полное правильное решение}
Пусть второй насос один заполняет резервуар за $x$ часов. Тогда за один час он заполняет $\frac{1}{x}$ резервуара. Первый насос за один час заполняет $\frac{1}{10}$ резервуара. Вместе они за один час заполняют $\frac{1}{6}$ резервуара. Поэтому
\[
\frac{1}{10}+\frac{1}{x}=\frac{1}{6}.
\]
Вычтем $\frac{1}{10}$ из обеих частей:
\[
\frac{1}{x}=\frac{1}{6}-\frac{1}{10}=\frac{1}{15}.
\]
Значит, $x=15$.

\subsection*{Проверка}
\[
\frac{1}{10}+\frac{1}{15}=\frac{3}{30}+\frac{2}{30}=\frac{5}{30}=\frac{1}{6}.
\]
Полученная общая производительность действительно соответствует заполнению резервуара за $6$ часов.

\subsection*{Итог}
\result{} второй насос один заполнит резервуар за $15$ часов.

\subsection*{Задача на закрепление}
Два насоса вместе заполняют резервуар за $4$ часа. Первый насос один заполняет его за $12$ часов. За сколько часов заполнит резервуар второй насос, работая один?

\newpage
\phantomsection\label{task:3}
\tasksection{3}
\subsection*{Текст задания}
Один насос заполняет резервуар на $4$ часа быстрее другого. Вместе они заполняют резервуар за $4$ часа $48$ минут. За сколько часов каждый насос заполняет резервуар отдельно?

\subsection*{Ваше решение}
Вы обозначили времена работы насосов через $x$ и $x+4$ и верно записали начальное уравнение
\[
\frac{1}{x+4}+\frac{1}{x}=\frac{1}{4{,}8}.
\]
Далее при избавлении от знаменателей появилась запись, приводящая к уравнению вида
\[
4{,}8x(x+4)=0,
\]
после чего было получено отрицательное значение для $x^2$, а затем оно заменено на положительное и получено $x=2$.

\subsection*{Ваш ответ}
В решении получено $2$ часа; два индивидуальных времени полностью не записаны.

\subsection*{Ошибка}
\err{} неверно выполнено избавление от знаменателей после правильно составленного уравнения.

\subsection*{Где именно возникает ошибка}
Последний правильный шаг:
\[
\frac{1}{x+4}+\frac{1}{x}=\frac{1}{4{,}8}.
\]
При приведении левой части к общему знаменателю должно появиться
\[
\frac{x+(x+4)}{x(x+4)}=\frac{2x+4}{x(x+4)},
\]
но числитель $2x+4$ в дальнейшем исчезает. Это и есть первый достоверно неверный переход.

\subsection*{Почему это неверно}
Сумма двух дробей с разными знаменателями не превращается в дробь с тем же знаменателем и нулевым числителем. Сначала нужно привести дроби к общему знаменателю $x(x+4)$. Числители при этом складываются: из первой дроби получается $x$, из второй --- $x+4$. Поэтому числитель равен $2x+4$.

Кроме того, из равенства $x^2=-4$ нельзя переходить к $x^2=4$ заменой отрицательного числа на модуль. Если квадрат действительного числа оказался равен отрицательному числу, это означает, что где-то ранее возникла ошибка или в действительных числах решений нет. В этой задаче положительное решение должно существовать, поэтому такой результат служит сигналом проверить предыдущие преобразования.

\subsection*{Ключевая идея}
\idea{} сначала складываем производительности, затем аккуратно приводим дроби к общему знаменателю:
\[
\frac{1}{x+4}+\frac{1}{x}=\frac{2x+4}{x(x+4)}.
\]

\subsection*{Исправление}
От последнего правильного шага:
\[
\frac{2x+4}{x(x+4)}=\frac{1}{4{,}8}.
\]
Умножаем обе части на $4{,}8x(x+4)$:
\[
4{,}8(2x+4)=x(x+4).
\]
Раскрываем скобки:
\[
9{,}6x+19{,}2=x^2+4x.
\]
Переносим всё в одну часть:
\[
x^2-5{,}6x-19{,}2=0.
\]
Чтобы убрать десятичные дроби, умножим на $5$:
\[
5x^2-28x-96=0.
\]
Дискриминант:
\[
D=(-28)^2-4\cdot5\cdot(-96)=784+1920=2704=52^2.
\]
Тогда
\[
x=\frac{28\pm52}{10}.
\]
Получаем $x=8$ или $x=-2{,}4$. Отрицательное время не подходит, поэтому $x=8$.
Второе время равно $x+4=12$.

\subsection*{Полное правильное решение}
Пусть более быстрый насос заполняет резервуар за $x$ часов. Тогда более медленный делает это за $x+4$ часа. Их производительности соответственно равны $\frac{1}{x}$ и $\frac{1}{x+4}$. Время совместной работы $4$ часа $48$ минут равно $4{,}8$ часа, поэтому общая производительность равна $\frac{1}{4{,}8}$.

Составим уравнение:
\[
\frac{1}{x}+\frac{1}{x+4}=\frac{1}{4{,}8}.
\]
После приведения к общему знаменателю и решения квадратного уравнения получаем единственное положительное значение $x=8$. Значит, второй насос работает отдельно $12$ часов.

\subsection*{Проверка}
\[
\frac{1}{8}+\frac{1}{12}=\frac{3}{24}+\frac{2}{24}=\frac{5}{24}.
\]
Обратная величина равна
\[
\frac{24}{5}=4{,}8 \text{ часа},
\]
то есть $4$ часа $48$ минут. Условие выполнено.

\subsection*{Итог}
\result{} насосы по отдельности заполняют резервуар за $8$ часов и за $12$ часов.

\subsection*{Задача на закрепление}
Один насос заполняет резервуар на $3$ часа быстрее другого. Вместе они заполняют резервуар за $3$ часа $36$ минут. За сколько часов каждый насос заполняет резервуар отдельно?

\newpage
\phantomsection\label{task:4}
\tasksection{4}
\subsection*{Текст задания}
Одна труба заполняет бассейн за $4$ часа, а другая труба полностью откачивает воду из заполненного бассейна за $12$ часов. За сколько часов наполнится пустой бассейн, если открыть обе трубы одновременно?

\subsection*{Ваше решение}
Записано
\[
\frac{1}{x}=\frac{1}{4}+\frac{1}{12},
\]
откуда получено $x=3$.

\subsection*{Ваш ответ}
$3$ часа.

\subsection*{Ошибка}
\err{} производительности труб сложены, хотя одна труба работает в противоположном направлении.

\subsection*{Где именно возникает ошибка}
Первый неверный шаг:
\[
\frac{1}{x}=\frac{1}{4}+\frac{1}{12}.
\]
Производительность откачивающей трубы должна вычитаться.

\subsection*{Почему это неверно}
Первая труба увеличивает количество воды в бассейне на $\frac14$ бассейна за час. Вторая труба, наоборот, уменьшает количество воды на $\frac1{12}$ бассейна за час. Если процессы направлены в разные стороны, их скорости изменения объёма имеют разные знаки. Поэтому итоговая скорость наполнения равна разности, а не сумме.

Сложение дало бы ситуацию, как будто обе трубы одновременно наливают воду. Но по условию вторая труба воду удаляет.

\subsection*{Ключевая идея}
\idea{} при одновременном наполнении и откачивании берём разность производительностей:
\[
\frac{1}{x}=\frac{1}{4}-\frac{1}{12}.
\]

\subsection*{Исправление}
\[
\frac{1}{x}=\frac{1}{4}-\frac{1}{12}=\frac{3}{12}-\frac{1}{12}=\frac{2}{12}=\frac{1}{6}.
\]
Следовательно,
\[
x=6.
\]

\subsection*{Полное правильное решение}
За один час первая труба наполняет $\frac14$ бассейна. За тот же час вторая труба откачивает $\frac1{12}$ бассейна. При одновременной работе чистый прирост воды за час равен
\[
\frac14-\frac1{12}=\frac16.
\]
Значит, за один час заполняется одна шестая бассейна, а весь бассейн наполнится за $6$ часов.

\subsection*{Проверка}
За $6$ часов первая труба подаст
\[
6\cdot\frac14=\frac32
\]
объёма бассейна, а вторая откачает
\[
6\cdot\frac1{12}=\frac12.
\]
Чистый объём:
\[
\frac32-\frac12=1.
\]
Получается ровно один полный бассейн.

\subsection*{Итог}
\result{} пустой бассейн наполнится за $6$ часов.

\subsection*{Задача на закрепление}
Одна труба заполняет бассейн за $5$ часов, а другая полностью откачивает воду из заполненного бассейна за $20$ часов. За сколько времени наполнится пустой бассейн при одновременной работе обеих труб?

\newpage
\phantomsection\label{task:5}
\tasksection{5}
\subsection*{Текст задания}
Для изготовления партии из $120$ деталей первый станок делает на $5$ деталей в час больше второго и поэтому заканчивает работу на $2$ часа раньше. Найдите производительность каждого станка.

\subsection*{Ваше решение}
Вы обозначили производительность второго станка через $x$, первого --- через $x+5$ и верно записали времена изготовления партии:
\[
\frac{120}{x}\quad\text{и}\quad\frac{120}{x+5}.
\]
Также верно составлено уравнение
\[
\frac{120}{x}-\frac{120}{x+5}=2
\]
и выполнен переход
\[
120(x+5)-120x=2x(x+5).
\]
После раскрытия скобок в следующем квадратном уравнении изменились знаки, и далее получен ответ
\[
\frac{-5+\sqrt{505}}{2}.
\]

\subsection*{Ваш ответ}
$\dfrac{-5+\sqrt{505}}{2}$ детали в час для одной из производительностей.

\subsection*{Ошибка}
\err{} при переносе членов после раскрытия скобок изменены знаки.

\subsection*{Где именно возникает ошибка}
Последний правильный шаг:
\[
120(x+5)-120x=2x(x+5).
\]
Раскрывая скобки, получаем
\[
120x+600-120x=2x^2+10x,
\]
то есть
\[
600=2x^2+10x.
\]
После переноса в левую часть должно быть
\[
2x^2+10x-600=0.
\]
В Вашем решении появляется уравнение с другими знаками, поэтому все дальнейшие вычисления относятся уже к другой задаче.

\subsection*{Почему это неверно}
При переносе слагаемого через знак равенства его знак меняется. Здесь $600$ переносится из левой части в правую или, эквивалентно, все члены собираются в одной части. Коэффициент при $10x$ не должен менять знак сам по себе: он уже находится в правой части вместе с $2x^2$. Ошибка в знаках меняет корни квадратного уравнения и, следовательно, производительности станков.

\subsection*{Ключевая идея}
\idea{} исходное уравнение составлено правильно; нужно сохранить его эквивалентность при раскрытии скобок и переносе членов.

\subsection*{Исправление}
Продолжаем от последнего правильного шага:
\[
600=2x^2+10x.
\]
Переносим $600$ вправо:
\[
2x^2+10x-600=0.
\]
Делим на $2$:
\[
x^2+5x-300=0.
\]
Дискриминант:
\[
D=5^2-4\cdot1\cdot(-300)=25+1200=1225=35^2.
\]
Тогда
\[
x=\frac{-5\pm35}{2}.
\]
Получаем $x=15$ или $x=-20$. Отрицательная производительность не подходит. Значит, второй станок делает $15$ деталей в час, а первый --- $20$ деталей в час.

\subsection*{Полное правильное решение}
Пусть второй станок производит $x$ деталей в час. Тогда первый производит $x+5$ деталей в час. На изготовление $120$ деталей второму требуется $\frac{120}{x}$ часов, а первому --- $\frac{120}{x+5}$ часов. Первый заканчивает на $2$ часа раньше, поэтому
\[
\frac{120}{x}-\frac{120}{x+5}=2.
\]
Умножаем на $x(x+5)$, где $x>0$:
\[
120(x+5)-120x=2x(x+5).
\]
Отсюда
\[
600=2x^2+10x,
\]
\[
x^2+5x-300=0.
\]
Положительный корень равен $15$. Следовательно, производительности равны $15$ и $20$ деталей в час.

\subsection*{Проверка}
Второй станок: $120:15=8$ часов. Первый станок: $120:20=6$ часов. Разница равна $2$ часам, как требуется. Кроме того, первый станок действительно производит на $5$ деталей в час больше.

\subsection*{Итог}
\result{} первый станок производит $20$ деталей в час, второй --- $15$ деталей в час.

\subsection*{Задача на закрепление}
Для изготовления партии из $96$ деталей первый станок делает на $4$ детали в час больше второго и заканчивает работу на $2$ часа раньше. Найдите производительность каждого станка.

\newpage
\phantomsection\label{task:6}
\tasksection{6}
\subsection*{Текст задания}
Две трубы вместе заполняют резервуар за $6$ часов. Первая труба в одиночку заполняет резервуар на $5$ часов быстрее второй. За сколько часов каждая труба заполнит резервуар отдельно?

\subsection*{Ваше решение}
Вы записали времена через $x$ и $x+5$ и верно начали с уравнения
\[
\frac{1}{x+5}+\frac{1}{x}=\frac{1}{6}.
\]
Далее дроби были преобразованы как будто знаменатели можно убрать без умножения всей равенства на общий знаменатель; появилась запись вида
\[
2x+5=\frac{1}{6},
\]
и в итоге получен ответ $\frac{29}{12}$.

\subsection*{Ваш ответ}
$\dfrac{29}{12}$ часа.

\subsection*{Ошибка}
\err{} при сложении и преобразовании дробей потерян общий знаменатель.

\subsection*{Где именно возникает ошибка}
Последний правильный шаг:
\[
\frac{1}{x+5}+\frac{1}{x}=\frac{1}{6}.
\]
Первый неверный переход --- замена суммы дробей выражением, в котором остаётся только числитель $2x+5$. На самом деле
\[
\frac{1}{x+5}+\frac{1}{x}=\frac{2x+5}{x(x+5)}.
\]

\subsection*{Почему это неверно}
При сложении дробей с разными знаменателями нельзя складывать только числители и забывать знаменатель. Общий знаменатель здесь $x(x+5)$. После приведения первой дроби к этому знаменателю её числитель становится $x$, второй --- $x+5$, поэтому сумма числителей действительно равна $2x+5$, но снизу обязательно остаётся $x(x+5)$.

Если этот знаменатель потерять, выражение меняет смысл: производительность, которая должна уменьшаться при увеличении времени, превращается в линейное выражение, которое, наоборот, растёт вместе с $x$.

\subsection*{Ключевая идея}
\idea{} после приведения к общему знаменателю получаем
\[
\frac{2x+5}{x(x+5)}=\frac{1}{6},
\]
и только затем избавляемся от знаменателей перекрёстным умножением.

\subsection*{Исправление}
\[
\frac{2x+5}{x(x+5)}=\frac{1}{6}.
\]
Отсюда
\[
6(2x+5)=x(x+5).
\]
Раскрываем скобки:
\[
12x+30=x^2+5x.
\]
Переносим всё в одну часть:
\[
x^2-7x-30=0.
\]
Разложим на множители:
\[
(x-10)(x+3)=0.
\]
Положительное решение $x=10$. Тогда второе время $x+5=15$.

\subsection*{Полное правильное решение}
Пусть первая, более быстрая труба заполняет резервуар за $x$ часов. Тогда вторая заполняет его за $x+5$ часов. Их производительности равны $\frac1x$ и $\frac1{x+5}$. Вместе они заполняют резервуар за $6$ часов, поэтому
\[
\frac1x+\frac1{x+5}=\frac16.
\]
После приведения к общему знаменателю и решения уравнения получаем $x=10$. Значит, первая труба заполняет резервуар за $10$ часов, а вторая --- за $15$ часов.

\subsection*{Проверка}
\[
\frac{1}{10}+\frac{1}{15}=\frac{3}{30}+\frac{2}{30}=\frac{5}{30}=\frac{1}{6}.
\]
Разница между индивидуальными временами равна $15-10=5$ часов. Оба условия выполнены.

\subsection*{Итог}
\result{} первая труба заполнит резервуар за $10$ часов, вторая --- за $15$ часов.

\subsection*{Задача на закрепление}
Две трубы вместе заполняют резервуар за $3$ часа $45$ минут. Первая труба в одиночку заполняет резервуар на $4$ часа быстрее второй. За сколько часов каждая труба заполнит резервуар отдельно?

\newpage
\phantomsection\label{task:7}
\tasksection{7}
\subsection*{Текст задания}
Две бригады одинаково квалифицированных рабочих одновременно начали выполнять два одинаковых заказа. В первой бригаде было $14$ рабочих, во второй --- $22$. Через $5$ дней из второй бригады в первую перешли $6$ рабочих. После перехода обе бригады закончили свои заказы одновременно. Сколько дней потребовалось на выполнение каждого заказа?

\subsection*{Ваше решение}
Вы правильно учли объём работы первых пяти дней: первая бригада выполняет $14\cdot5=70$ условных единиц, вторая --- $22\cdot5=110$. После перехода рабочих в первой бригаде стало $20$ человек, во второй --- $16$. Для числа дней после перехода записано равенство
\[
70+20y=110+16y.
\]
Из него верно получено $y=10$. Затем записан ответ $10$ дней.

\subsection*{Ваш ответ}
$10$ дней.

\subsection*{Ошибка}
\err{} найдено число дней после перехода рабочих, но оно принято за полное время выполнения заказа.

\subsection*{Где именно возникает ошибка}
Последний правильный шаг:
\[
y=10.
\]
Это означает: после перехода шести рабочих обеим бригадам понадобилось ещё $10$ дней. Первый неверный шаг --- записать $10$ как ответ на вопрос о полном времени работы, не прибавив уже прошедшие $5$ дней.

\subsection*{Почему это неверно}
Переменная $y$ введена для второго этапа работы, который начинается только после первых пяти дней. В условии спрашивается, сколько дней заняло выполнение заказа от самого начала до завершения. Поэтому полное время состоит из двух последовательных промежутков: первых $5$ дней и последующих $10$ дней.

Алгебраическая часть решения у Вас здесь фактически верна; ошибка возникает при интерпретации найденной величины.

\subsection*{Ключевая идея}
\idea{} всегда проверяйте, что именно обозначает найденная переменная: здесь $y$ --- только длительность второго этапа, а не всего процесса.

\subsection*{Исправление}
После полученного Вами $y=10$ нужно дописать:
\[
5+10=15.
\]

\subsection*{Полное правильное решение}
Пусть производительность одного рабочего за день равна одной условной единице; конкретное значение не важно, потому что все рабочие одинаково квалифицированы, а заказы одинаковы.

За первые $5$ дней первая бригада выполнит $14\cdot5=70$ условных единиц, вторая --- $22\cdot5=110$.

После перехода $6$ рабочих численность бригад станет $20$ и $16$. Пусть после перехода они работают ещё $y$ дней. Поскольку заказы одинаковы и заканчиваются одновременно, их полные объёмы равны:
\[
70+20y=110+16y.
\]
Отсюда
\[
4y=40,
\]
\[
y=10.
\]
Это только время после перехода. Полное время:
\[
5+10=15 \text{ дней}.
\]

\subsection*{Проверка}
Первая бригада выполнит
\[
14\cdot5+20\cdot10=70+200=270
\]
условных единиц. Вторая:
\[
22\cdot5+16\cdot10=110+160=270.
\]
Объёмы равны, значит одинаковые заказы действительно завершены одновременно.

\subsection*{Итог}
\result{} каждый заказ выполнялся $15$ дней.

\subsection*{Задача на закрепление}
Две бригады одинаково квалифицированных рабочих одновременно начали выполнять два одинаковых заказа. В первой бригаде было $12$ рабочих, во второй --- $20$. Через $4$ дня из второй бригады в первую перешли $6$ рабочих. После перехода обе бригады закончили свои заказы одновременно. Сколько дней потребовалось на выполнение каждого заказа?

\newpage
\phantomsection\label{task:8}
\tasksection{8}
\subsection*{Текст задания}
Анна и Борис вместе выполняют работу за $8$ часов, Борис и Виктор --- за $12$ часов, Виктор и Анна --- за $8$ часов. За сколько часов они выполнят эту работу втроём?

\subsection*{Ваше решение}
В таблице записаны времена парной работы $8$, $12$ и $8$ часов. Затем эти времена складываются, получается $28$, после чего сумма делится на $2$ и получается $14$. Отдельная строка с окончательным ответом не записана, но видимый вывод вычислений равен $14$.

\subsection*{Ваш ответ}
Видимый итог вычислений --- $14$ часов.

\subsection*{Ошибка}
\err{} складывались времена выполнения работы, хотя складывать нужно производительности.

\subsection*{Где именно возникает ошибка}
Первый неверный шаг --- использование суммы
\[
8+12+8=28
\]
как величины, из которой можно получить время общей работы втроём. Часы парной работы не являются долями одной и той же величины, которые можно складывать таким способом.

\subsection*{Почему это неверно}
Если пара выполняет всю работу за $8$ часов, это означает, что за один час эта пара делает $\frac18$ работы. Именно $\frac18$ является производительностью пары. Аналогично для времени $12$ часов производительность равна $\frac1{12}$.

Когда складываются уравнения для трёх пар, каждый человек оказывается учтён дважды. Поэтому сумма трёх парных производительностей равна удвоенной производительности всей тройки. Это позволяет быстро найти общую производительность.

\subsection*{Ключевая идея}
\idea{} переводим каждое время в производительность, то есть берём обратную величину.

\subsection*{Исправление}
Обозначим индивидуальные производительности через $a$, $b$, $c$. Тогда
\[
a+b=\frac18,
\]
\[
b+c=\frac1{12},
\]
\[
c+a=\frac18.
\]
Складываем:
\[
2(a+b+c)=\frac18+\frac1{12}+\frac18.
\]
Правая часть:
\[
\frac18+\frac18+\frac1{12}=\frac14+\frac1{12}=\frac{3}{12}+\frac1{12}=\frac13.
\]
Значит,
\[
a+b+c=\frac16.
\]
Следовательно, втроём они выполнят работу за $6$ часов.

\subsection*{Полное правильное решение}
Пусть $a$, $b$, $c$ --- части работы, которые Анна, Борис и Виктор выполняют за один час. По условию:
\[
a+b=\frac18,\qquad b+c=\frac1{12},\qquad c+a=\frac18.
\]
Сумма трёх равенств даёт
\[
2a+2b+2c=\frac13.
\]
Делим на $2$:
\[
a+b+c=\frac16.
\]
Это означает, что втроём за один час они выполняют одну шестую всей работы. Следовательно, на всю работу требуется $6$ часов.

\subsection*{Проверка}
Общая производительность $\frac16$ означает, что за $6$ часов будет выполнено
\[
6\cdot\frac16=1
\]
полная работа.

\subsection*{Итог}
\result{} втроём Анна, Борис и Виктор выполнят работу за $6$ часов.

\subsection*{Задача на закрепление}
Первый и второй исполнители вместе выполняют работу за $6$ часов, второй и третий --- за $8$ часов, третий и первый --- за $12$ часов. За сколько времени они выполнят эту работу втроём?

\newpage
\phantomsection\label{task:9}
\tasksection{9}
\subsection*{Текст задания}
Первый и второй мастера вместе выполняют работу за $6$ часов, второй и третий --- за $10$ часов, третий и первый --- за $7{,}5$ часа. За сколько часов выполнил бы эту работу каждый мастер отдельно?

\subsection*{Ваше решение}
В таблице записаны времена парной работы $6$, $10$ и $7{,}5$ часа. Далее эти времена складываются и делятся на $2$, получается $11{,}75$, после чего из этой величины вычитаются парные времена. В правой таблице записаны индивидуальные времена $1{,}75$, $4{,}25$ и $5{,}25$ часа.

\subsection*{Ваш ответ}
В таблице указаны $1{,}75$, $4{,}25$ и $5{,}25$ часа.

\subsection*{Ошибка}
\err{} времена парной работы обрабатываются как обычные слагаемые; нужно сначала перейти к производительностям.

\subsection*{Где именно возникает ошибка}
Первый неверный шаг --- вычисление
\[
\frac{6+10+7{,}5}{2}=11{,}75
\]
с целью получить суммарную характеристику работы мастеров. Время выполнения работы не складывается линейно. Складываются доли работы, выполняемые за один час.

\subsection*{Почему это неверно}
Пара, работающая $6$ часов, имеет производительность $\frac16$ работы в час; пара, работающая $10$ часов, --- $\frac1{10}$; пара, работающая $7{,}5$ часа, --- $\frac1{7{,}5}=\frac2{15}$. Только эти величины можно складывать и вычитать как скорости выполнения одной и той же работы.

После нахождения общей производительности всех трёх мастеров индивидуальную производительность каждого удобно получить вычитанием производительности противоположной пары. Например, производительность первого мастера равна общей производительности всех троих минус производительность второго и третьего вместе.

\subsection*{Ключевая идея}
\idea{} решаем задачу в производительностях, а индивидуальное время получаем в самом конце как обратную величину.

\subsection*{Исправление}
Обозначим индивидуальные производительности через $a$, $b$, $c$. Тогда
\[
a+b=\frac16,
\]
\[
b+c=\frac1{10},
\]
\[
c+a=\frac1{7{,}5}=\frac2{15}.
\]
Складываем:
\[
2(a+b+c)=\frac16+\frac1{10}+\frac2{15}.
\]
Приводим к знаменателю $30$:
\[
2(a+b+c)=\frac5{30}+\frac3{30}+\frac4{30}=\frac{12}{30}=\frac25.
\]
Значит,
\[
a+b+c=\frac15.
\]
Теперь
\[
a=\frac15-\frac1{10}=\frac1{10},
\]
\[
b=\frac15-\frac2{15}=\frac1{15},
\]
\[
c=\frac15-\frac16=\frac1{30}.
\]
Следовательно, индивидуальные времена равны $10$, $15$ и $30$ часов.

\subsection*{Полное правильное решение}
Пусть первый мастер за час выполняет $a$ часть работы, второй --- $b$, третий --- $c$. Тогда из условий о парной работе:
\[
a+b=\frac16,\qquad b+c=\frac1{10},\qquad c+a=\frac2{15}.
\]
Сложив все три равенства, получаем
\[
2(a+b+c)=\frac25,
\]
поэтому
\[
a+b+c=\frac15.
\]
Производительность первого мастера:
\[
a=(a+b+c)-(b+c)=\frac15-\frac1{10}=\frac1{10}.
\]
Значит, первому одному нужно $10$ часов.

Производительность второго:
\[
b=(a+b+c)-(a+c)=\frac15-\frac2{15}=\frac1{15},
\]
поэтому второму одному нужно $15$ часов.

Производительность третьего:
\[
c=(a+b+c)-(a+b)=\frac15-\frac16=\frac1{30},
\]
поэтому третьему одному нужно $30$ часов.

\subsection*{Проверка}
Первый и второй:
\[
\frac1{10}+\frac1{15}=\frac16,
\]
то есть вместе $6$ часов.

Второй и третий:
\[
\frac1{15}+\frac1{30}=\frac1{10},
\]
то есть вместе $10$ часов.

Третий и первый:
\[
\frac1{30}+\frac1{10}=\frac4{30}=\frac2{15}=\frac1{7{,}5},
\]
то есть вместе $7{,}5$ часа. Все условия выполняются.

\subsection*{Итог}
\result{} первый мастер выполнит работу один за $10$ часов, второй --- за $15$ часов, третий --- за $30$ часов.

\subsection*{Задача на закрепление}
Первый и второй мастера вместе выполняют работу за $4$ часа $48$ минут, второй и третий --- за $8$ часов, третий и первый --- за $6$ часов. За сколько часов выполнил бы эту работу каждый мастер отдельно?

\vfill
\begin{center}
\small Лёвин Артём Александрович эксклюзивно для Анастасии Павловой
\end{center}

\end{document}
